% Pursuing Stacks (cotes 134-2 to 134-8): archive cote 134-6, site page 52 (PDF page 53), typescript page 370.
% typescript. Draft compared: ps5.tex:3898-4016.
% AI-transcribed by Claude Opus 5.5 (claude-opus-5-5, adaptive thinking, effort high), Message Batches API, batch ps3 (msgbatch_01XdFmuVVBAcgc7EFxdtCqsj), 2026-11-12; not yet checked by eye.
\scanpage{134-6/52}
\mspage{370}

I played around some yesterday and today with the formalism of integra-%
tors,  notably with respect to maps
\[ f : B' \longrightarrow B \]
between small categories,  and the \inserted{corresponding} \uline{integration functor}
\[ f_{!}^{ab} : B'^{\wedge}_{ab} \longrightarrow B^{\wedge}_{ab} \quad , \]
and it's left derived functor $Lf_{!}^{ab}$ .  Thus the chain complex in $B^{\wedge}_{ab}$
\margin{$L^{f}_{\bullet}$} \hand{(14)} \[ L_{\bullet}^{B'/B} \quad \text{or} \quad L_{\bullet}^{f} \overset{\text{def}}{=} f_{!}^{ab}(L_{\bullet}^{B'}) \quad , \]
which \xout{is a projective}  has projective components (and even is \unsure{a\xout{n}} chain
complex in Addinf(B)  resp. in Add(B),  if $L_{\bullet}^{B'}$ is quasi-spe\xout{\illegible{1}}cial  resp. is
special),  and is defined up to chain homoto\unsure{\hand{p}}ism,  embodies the relative
homology properties of B' over B  i.e. of f ,  in much the same way as (11)
embodies the glo\unsure{\xout{v}}bal homology properties of B  (i.e. of B over one point).
When the functor f is ``coaspheric''  i.e. the functor
\[ f^{o} : B'^{o} \longrightarrow B^{o} \]
between the opposite categories is aspheric,  then $L_{\bullet}^{B'/B}$ is again an inte-%
grator on B,  and the converse should hold too provided we take the meaning
of ``coaspheric'' and ``aspheric'' with respect to a suitable basic localizer
$\underline{W} = \underline{W}_{\infty}^{Z}$  -  presumably we'll come back upon this in part V or part VI of
the notes.  For the time being,  it seems more interesting to give now the
precise relationship between the notion of an \uline{integrator}\xout{,} for B,  and the
notion of an \uline{abelianizator} for the dual category $A = B^{o}$ ,  introduced in
section 93.

\uline{Remarks}  1)  It is a familiar fact that when working in \hand{\v{}}Cech-flavored
contexts,  such as general topoi,  or etale topoi for schemes and the like,
one has throughout and from the start a good hold upon \uline{cohomology} notions,
whereas it is a lot more subtle to squeeze out adequate homology notions,
which (to my knowledge) can be carried through only indirectly via cohomolo%
-gy,  and using suitable finiteness and duality statements within the cohomo%
-logy formalism.  Historically however,  homology was introduced before
cohomology via cellular decompositions of spaces,  with a more direct ap-%
peal to geometric intuition.  This preference for homology rather than co-%
homology seems to be \xout{a} still prevalent among most homotopy theorists, who
have a tendency to view a topological space (however wild it may be) as
being no more no less than it's singular complex.  A comprehensive state-%
ment establishing,  in a suitable wide enough context,  essential equiva-%
lence between the two viewpoints,  seems to be still lacking,  as far as I

\margin{52}
